\begin{lemma}[$L^2(\gamma_n)$ lies in $L^1_{\mathrm{loc}}$ and in $L^2_{\mathrm{loc}}$ for Lebesgue measure]
Let $\gamma_n$ be as in \noderef{GaussMeasure}, write $\lambda$ for Lebesgue
measure on $\mathbb{R}^n$, and let $v : \mathbb{R}^n \to \mathbb{R}$ be
$\gamma_n$-measurable with $\int_{\mathbb{R}^n} |v|^2 \, d\gamma_n < \infty$,
i.e.\ $v \in L^2(\gamma_n)$ as a function.  Then:
\begin{enumerate}
\item $v$ is locally $\lambda$-integrable, i.e.\ every point of $\mathbb{R}^n$ has
a neighbourhood $U$ with $\int_U |v| \, d\lambda < \infty$;
\item for every compact $K \subseteq \mathbb{R}^n$ the restriction of $v$ to $K$
lies in $L^2(K,\lambda)$, i.e.\ $\int_K |v|^2 \, d\lambda < \infty$;
\item (quantitative form, uniform in the function) for every compact
$K \subseteq \mathbb{R}^n$ there is a constant $C_K \geq 0$, depending only on
$K$ and on $n$, such that \emph{every} $\gamma_n$-measurable
$g : \mathbb{R}^n \to \mathbb{R}$ with
$\int_{\mathbb{R}^n} |g|^2 \, d\gamma_n < \infty$ is $\lambda$-integrable on $K$
and satisfies
\[
  \int_K |g| \, d\lambda \;\leq\; C_K\, \|g\|_{L^2(\gamma_n)} .
\]
\end{enumerate}
All three conclusions are statements about Lebesgue measure, not about
$\gamma_n$; they are the comparison that lets a function of the Gaussian space
be fed to test identities written against Lebesgue measure.  Conclusion (3)
differs from (1) in that its constant does not depend on the function, which is
what permits passing to a limit along an $L^2(\gamma_n)$-convergent sequence
inside an integral taken against $\lambda$.
\end{lemma}
\begin{proof}
\emph{Step 0: measurability transfers between the two measures.}
By \noderef{GaussMeasure}, $\gamma_n$ is absolutely continuous with respect to
$\lambda$ with Radon--Nikodym density $\varrho(x) = \pi^{-n/2}e^{-|x|^2}$, which
is continuous and strictly positive at every point; in particular
$\int h \, d\gamma_n = \int h\,\varrho \, d\lambda$ for every measurable
$h : \mathbb{R}^n \to [0,+\infty]$, the case $h = \mathbf{1}_A$ being the
set-level formula displayed there.  Consequently a set is $\lambda$-null if and
only if it is $\gamma_n$-null: if $\lambda(A) = 0$ then
$\gamma_n(A) = \int_A \varrho \, d\lambda = 0$; and if $\gamma_n(A) = 0$ then,
since $\varrho > 0$ everywhere, $\int_A \varrho\,d\lambda = 0$ forces
$\lambda(A) = 0$.  Hence $v$, which is $\gamma_n$-almost everywhere equal to a
measurable function, is also $\lambda$-almost everywhere equal to that same
function, so $v$ is $\lambda$-measurable.  All integrals below are therefore
meaningful.

\emph{Step 1: the density comparison on a compact set.}
Let $K \subseteq \mathbb{R}^n$ be compact and nonempty (for $K = \emptyset$ both
conclusions are trivial).  Since $\varrho$ is continuous and $K$ is compact and
nonempty, $\varrho$ attains a minimum on $K$ by the extreme value theorem, which
is shared context; call it $c_K = \min_{x \in K} \varrho(x) = \varrho(x_-)$ for a
minimizing $x_- \in K$.  As $\varrho > 0$ everywhere, $c_K > 0$.  Also
$\lambda(K) < \infty$, again shared context for compact sets.  Now
let $h : \mathbb{R}^n \to [0,+\infty]$ be measurable.  On $K$ we have
$h \leq c_K^{-1}\, h\, \varrho$ pointwise, because $\varrho \geq c_K$ there and
$h \geq 0$.  Integrating this pointwise inequality over $K$ and then using the
density formula of \noderef{GaussMeasure},
\begin{equation*}
  \int_K h \, d\lambda
    \;\leq\; \frac{1}{c_K} \int_K h\, \varrho \, d\lambda
    \;=\; \frac{1}{c_K} \int_K h \, d\gamma_n .
  \tag{1}
\end{equation*}

\emph{Step 2: conclusion (2).}
Apply $(1)$ with $h = |v|^2$:
\[
  \int_K |v|^2 \, d\lambda
    \;\leq\; \frac{1}{c_K} \int_K |v|^2 \, d\gamma_n
    \;\leq\; \frac{1}{c_K} \int_{\mathbb{R}^n} |v|^2 \, d\gamma_n
    \;=\; \frac{1}{c_K}\, \|v\|_{L^2(\gamma_n)}^2
    \;<\; \infty ,
\]
the middle inequality because $|v|^2 \geq 0$, and finiteness by the hypothesis
$\int |v|^2 d\gamma_n < \infty$.  So $v \in L^2(K,\lambda)$, which is conclusion (2).

\emph{Step 3: conclusion (1).}
Fix $x_0 \in \mathbb{R}^n$ and put $K = \overline{B}(x_0,1)$, a compact
neighbourhood of $x_0$; $\lambda(K) < \infty$ by the shared-context item on
compact sets.  By Cauchy--Schwarz, which is shared context, applied on the
measure space $(K,\lambda)$ to the pair of functions $|v|$ and the constant
function $1$ --- both in $L^2(K,\lambda)$, the first by Step 2 and the second
because $\lambda(K) < \infty$ ---
\[
  \int_K |v| \, d\lambda
    \;\leq\; \Bigl(\int_K |v|^2 \, d\lambda\Bigr)^{1/2}
             \Bigl(\int_K 1 \, d\lambda\Bigr)^{1/2}
    \;=\; \Bigl(\int_K |v|^2 \, d\lambda\Bigr)^{1/2} \lambda(K)^{1/2} ,
\]
which is finite by Step 2 and $\lambda(K) < \infty$.  So $v$ is
$\lambda$-integrable on the neighbourhood $K$ of $x_0$.  As $x_0$ was arbitrary,
$v$ is locally $\lambda$-integrable, which is conclusion (1).

\emph{Step 4: conclusion (3), with an explicit constant.}
Let $K$ be compact.  If $K$ is empty put $C_K = 0$; the asserted inequality then
reads $0 \leq 0$.  Otherwise set
\[
  C_K := \Bigl(\frac{\lambda(K)}{c_K}\Bigr)^{1/2},
\]
with $c_K = \min_{x\in K}\varrho(x) > 0$ as in Step 1; this is a finite
nonnegative number depending only on $K$ and on $n$, because $K$ compact gives
$\lambda(K) < \infty$ and Step 1 gives $c_K > 0$.  Let $g$ be any
$\gamma_n$-measurable function with $\int |g|^2 \, d\gamma_n < \infty$.  Steps 0
and 1 apply verbatim with $g$ in place of $v$ --- Step 0 used only
$\gamma_n$-measurability of $v$, and Step 1 is a statement about an arbitrary
nonnegative measurable $h$ and does not mention $v$ at all --- so $g$ is
$\lambda$-measurable and, by $(1)$ with $h = |g|^2$ together with
$|g|^2 \geq 0$,
\[
  \int_K |g|^2 \, d\lambda
    \;\leq\; \frac{1}{c_K}\int_K |g|^2\, d\gamma_n
    \;\leq\; \frac{1}{c_K}\int_{\mathbb{R}^n} |g|^2\, d\gamma_n
    \;=\; \frac{1}{c_K}\,\|g\|_{L^2(\gamma_n)}^2 .
\]
Taking square roots, which preserves the inequality because $t \mapsto t^{1/2}$
is nondecreasing on $[0,+\infty)$ and both sides are nonnegative and finite,
gives $\bigl(\int_K|g|^2 d\lambda\bigr)^{1/2} \leq c_K^{-1/2}
\|g\|_{L^2(\gamma_n)}$.  By Cauchy--Schwarz, which is shared context, on the
measure space $(K,\lambda)$ applied to $|g|$ and the constant function $1$ ---
both in $L^2(K,\lambda)$, the first by the display just proved and the second
because $\lambda(K) < \infty$ ---
\[
  \int_K |g| \, d\lambda
    \;\leq\; \Bigl(\int_K |g|^2 \, d\lambda\Bigr)^{1/2} \lambda(K)^{1/2}
    \;\leq\; \frac{\lambda(K)^{1/2}}{c_K^{1/2}}\,\|g\|_{L^2(\gamma_n)}
    \;=\; C_K\, \|g\|_{L^2(\gamma_n)} ,
\]
the second step multiplying the previous bound by $\lambda(K)^{1/2} \geq 0$.
Finally, $g$ is $\lambda$-integrable on $K$: it is $\lambda$-measurable by Step 0,
and the Cauchy--Schwarz bound just displayed is precisely the finiteness of
$\int_K |g|\, d\lambda$, so $g$ restricted to $K$ lies in $L^1(K,\lambda)$.
Since $g$ was arbitrary, this is conclusion (3).
\end{proof}
